% answers_short.tex —— 2006 年·普通高等学校招生考试·江西卷文科 · 简短答案速查
% 由 tools/make_short_answers.py 自 answers.tex 的 \gksolhead/\gkanswer 提取，
% 并与 2006/2006_数学真题_江西文_答案要点.md 逐题比对后生成，勿手改（改源 MD 后重跑）。

\begin{gknotice}
\textbf{使用说明}：本节只给最终答案便于快速核对，完整推导见解析版。答案为编写者独立推导，\textbf{非官方标准答案}。
\end{gknotice}

\begin{gknotice}
选择题答案中的 U 表示待核，未给出猜测答案，详情见解析版。
\end{gknotice}

\gkpart{选择题}
\begin{center}\begin{tabular}{ccccccccccccc}
\toprule
题号 & 1 & 2 & 3 & 4 & 5 & 6 & 7 & 8 & 9 & 10 & 11 & 12 \\
\midrule
答案 & C & B & A & U & C & C & B & U & B & A & D & D \\
\bottomrule
\end{tabular}\end{center}

\gkpart{填空题}
\begin{itemize}[leftmargin=2.6em]
  \item[\textbf{13.}] $\sqrt2$
  \item[\textbf{14.}] $2$
  \item[\textbf{15.}] $10$
  \item[\textbf{16.}] ①④
\end{itemize}

\gkpart{解答题}
\begin{itemize}[leftmargin=2.6em]
  \item[\textbf{17.}] (1) $a=-\dfrac12\text{，}b=-2$；单调递增区间为 $(-\infty,-\frac23)$、$(1,+\infty)$，单调递减区间为 $(-\frac23,1)$；(2) $c\in(-\infty,-1)\cup(2,+\infty)$
  \item[\textbf{18.}] (1) $\dfrac{729}{1000}$；(2) $\dfrac{131}{500}$
  \item[\textbf{19.}] (1) $\dfrac73$；(2) $b=\sqrt3$
  \item[\textbf{20.}] (1) $\sqrt{\dfrac23}$；(2) $\arccos\dfrac25$；(3) $\arccos\dfrac7{3\sqrt6}$
  \item[\textbf{21.}] (1) $\dfrac{(x-c/2)^2}{a^2}+\dfrac{y^2}{b^2}=\dfrac{c^2}{4a^2}$；(2) $\theta=2\arctan\dfrac12$（此时 $\sin\theta=\dfrac45\text{，}\cos\theta=\dfrac35$）
  \item[\textbf{22.}] (1) 令 $x_n=a_n-\dfrac1{a_n}$，则 $x_n=\dfrac83\,2^{n-1}$，因此 $a_n=\frac{x_n+\sqrt{x_n^2+4}}2$；(2) $S_n+T_n=\dfrac{64(4^n-1)}{27}+2n$；使其为整数的最小正整数为 $n=9$
\end{itemize}
